\documentclass[12pt,a4paper]{article}
\usepackage[utf8]{inputenc}
\usepackage[T1]{fontenc}
\usepackage{amsmath,amssymb,amsthm,mathrsfs,mathtools}
\usepackage[margin=1in]{geometry}
\usepackage{hyperref}
\hypersetup{colorlinks=true,linkcolor=blue,citecolor=blue,urlcolor=blue}
\usepackage{booktabs}

\theoremstyle{plain}
\newtheorem{theorem}{Theorem}[section]
\newtheorem{proposition}[theorem]{Proposition}
\newtheorem{lemma}[theorem]{Lemma}
\newtheorem{corollary}[theorem]{Corollary}
\theoremstyle{definition}
\newtheorem{definition}[theorem]{Definition}
\theoremstyle{remark}
\newtheorem{remark}[theorem]{Remark}

\newcommand{\Disc}{\operatorname{Disc}}
\newcommand{\dd}{\mathrm{d}}
\newcommand{\R}{\mathbb{R}}
\newcommand{\C}{\mathbb{C}}
\newcommand{\cM}{\mathcal{M}}
\newcommand{\dLIPS}{\dd\Pi_2}

\title{Loop Amplitudes from Unitarity Cuts in Celestial Mellin Space:\\[4pt]
A Complete Derivation for the Scalar Box}

\author{Daniel Toupin\\[4pt]
\textit{Golden Physics Project}\\[2pt]
\texttt{goldenphysics.org}}

\date{August 2026 --- v19, superseding arXiv-style drafts v14--v18 and PhilArchive TOULAF-2}

\begin{document}
\maketitle

\begin{abstract}
\noindent
We derive, with every Jacobian computed and every constant fixed by
first principles, the complete chain that turns tree-level data into a
one-loop amplitude in celestial (Mellin) space.  The chain has four
links, each verified numerically to between $7$ and $30$ significant
digits with \emph{zero} fitted parameters.
(i)~The two-particle massless unitarity cut, written in celestial
coordinates, is uniform on the round celestial sphere with measure
$\dd^2z/\bigl(8\pi^2(1+|z|^2)^2\bigr)$, and momentum conservation forces
the cut pair to antipodal points $z_6=-1/\bar z_5$.
(ii)~The Mellin image of the cut measure is
$\Phi(\Delta_5,\Delta_6;M)=\tfrac{1}{8\pi}(M/2)^{\Delta_5+\Delta_6-2}
\,\Gamma(\Delta_5)\Gamma(\Delta_6)/\Gamma(\Delta_5+\Delta_6)$: the
celestial three-point (OPE) Gamma structure is a theorem of phase space,
and the shadow locus $\Delta_5+\Delta_6=2$ is the locus of exact scale
invariance of the cut, not a pole at which a residue is taken.
(iii)~For the internally regulated scalar box the $s$-channel cut
reduces to a two-dimensional celestial integral with the exact closed
form $C=\tfrac{1}{8\pi}\int_0^1\dd x\,[\mu^2(s+\mu^2)+s(s+t)x(1-x)]^{-1}$;
its Mellin transform at fixed angle has a universal, infrared-safe
double pole $\tfrac{1}{4\pi\kappa}(2-\sigma)^{-2}$, with the collinear
logarithm confined to the subleading simple pole.
(iv)~The full loop is recovered from the cut by an unsubtracted
fixed-$u$ dispersion relation whose normalization,
$\operatorname{Im}J=8\pi^2 C$, is derived and anchored exactly on the
bubble; the reconstruction agrees with the directly computed box to
machine precision ($\lesssim 10^{-21}$) at all tested kinematics.  In
Mellin space the entire dispersion collapses to multiplication by a
universal kernel,
\[
\cM_{J}(\sigma)\;=\;\frac{8\pi^2}{\sin(\pi\sigma)}\;\cM_{C}(\sigma),
\qquad 0<\operatorname{Re}\sigma<1 ,
\]
which we propose as the correct, derived form of the ``loop integrands
from shadow discontinuities'' program.  The spectral weight
$P(\lambda)=\pi\lambda/\sinh(\pi\lambda)$ used in earlier drafts is
recovered here as the \emph{derived} phase-space factor
$\Gamma(1{+}i\lambda)\Gamma(1{-}i\lambda)$ on the shadow locus; it is
not a replacement for the Feynman measure and it does not by itself
render loop expansions ultraviolet finite.  We state precisely which
claims of the earlier drafts survive, which are corrected, and which
are withdrawn.
\end{abstract}

\medskip
\noindent\textbf{Keywords:} celestial holography, unitarity, Cutkosky
rules, Mellin amplitudes, dispersion relations, scalar box.

\tableofcontents
\newpage

% =====================================================================
\section{Introduction and statement of results}
% =====================================================================

Generalized unitarity reconstructs loop amplitudes from products of
trees sewn across cuts \cite{Bern1994}.  Celestial holography rewrites
scattering in a boost (Mellin) basis on the celestial sphere
\cite{Pasterski2017,Raclariu2021}.  Earlier drafts in this series
(v14--v18, and the PhilArchive record TOULAF-2) proposed that loop
integrands are encoded in tree celestial amplitudes and extracted by
``shadow discontinuities'' supported on the locus
$\Delta_5+\Delta_6=2$.  The proposal contained a true statement wrapped
in an incorrect mechanism, and its numerical support did not withstand
audit: the archived benchmark \texttt{check\_box} compared a
Feynman-parameter integral against a reparameterization of itself, and
the published closed-form checks of the v3 lineage fail recomputation.

This paper replaces the mechanism with a derivation.  The strategy is
the only unconditionally safe one available: we do not guess a
celestial extraction rule and match it to a known answer; we
Mellin-transform the one true loops-from-trees statement in quantum
field theory --- the Cutkosky unitarity relation --- computing every
Jacobian along the way.  What emerges is a four-link chain, each link a
theorem verified numerically with no adjustable constants:

\begin{enumerate}
\item[\textbf{L1.}] \emph{Cut geometry} (Thm.~\ref{thm:measure},
\ref{thm:phasespace}): the two-particle massless phase space in
celestial coordinates, including the antipodal pairing of the cut legs.
\item[\textbf{L2.}] \emph{Mellin image of the cut}
(Thm.~\ref{thm:mellincut}): the
$\Gamma\Gamma/\Gamma$ structure and the true meaning of the shadow
locus.
\item[\textbf{L3.}] \emph{The box cut and its Mellin poles}
(Thm.~\ref{thm:boxcut}, \ref{thm:poles}): an exact closed form and the
universal infrared-safe double pole.
\item[\textbf{L4.}] \emph{Reconstruction} (Thm.~\ref{thm:disp},
\ref{thm:celdisp}): the loop from its cut by dispersion, and its
Mellin-space form as multiplication by $8\pi^2/\sin(\pi\sigma)$.
\end{enumerate}

Section~\ref{sec:relation} states the corrected relation to the earlier
shadow-discontinuity program and to the companion drafts on the
``Plancherel loop measure,'' the kinematic block, and the blackbody
analogy.  Section~\ref{sec:notclaimed} lists what is \emph{not} claimed.

% =====================================================================
\section{Conventions}
\label{sec:conventions}
% =====================================================================

Metric $(+,-,-,-)$.  A massless momentum is parameterized as
\begin{equation}
\ell \;=\; \omega\, q(z,\bar z),\qquad
q(z,\bar z) \;=\; \bigl(1+|z|^2,\; z+\bar z,\; -i(z-\bar z),\;
1-|z|^2\bigr),\qquad \omega>0,
\label{eq:param}
\end{equation}
so that $q^2=0$ identically and, for two celestial points,
\begin{equation}
q(z)\cdot q(w) \;=\; 2\,|z-w|^2 .
\label{eq:qq}
\end{equation}
Loop integrals are normalized as $I=\int
\frac{\dd^4\ell}{(2\pi)^4}\prod_i (D_i+i\varepsilon)^{-1}$.  Two-particle
Lorentz-invariant phase space is
$\dLIPS=(2\pi)^4\delta^4(P-\ell_5-\ell_6)\prod_{i=5,6}
\frac{\dd^3\ell_i}{(2\pi)^3 2E_i}$, with the standard massless total
$\int\dLIPS=\tfrac{1}{8\pi}$.  All numerical statements below are
reproduced by the scripts listed in Sec.~\ref{sec:code}; quoted
relative errors are those of the archived runs.

% =====================================================================
\section{Link 1: the unitarity cut on the celestial sphere}
\label{sec:link1}
% =====================================================================

\begin{lemma}[On-shell measure]
\label{lem:onshell}
With the parameterization \eqref{eq:param},
$\displaystyle \frac{\dd^3\ell}{2E} = 2\,\omega\,\dd\omega\,\dd^2z$,
where $\dd^2z=\dd x\,\dd y$.  The Jacobian
$\det \partial(\ell^1,\ell^2,\ell^3)/\partial(\omega,x,y)
=4\omega^2(1+|z|^2)$ is computed symbolically in the companion script.
\end{lemma}

\begin{theorem}[Cut geometry: antipodal pairing and uniform measure]
\label{thm:measure}
Let $P=(M,0,0,0)$, $M>0$.  The constraint
$\delta^4(P-\ell_5-\ell_6)$ with both legs massless and forward is
solved \emph{uniquely} by
\begin{equation}
z_6=-\frac{1}{\bar z_5},\qquad
\omega_5=\frac{M}{2(1+|z_5|^2)},\qquad
\omega_6=\frac{M\,|z_5|^2}{2(1+|z_5|^2)},
\label{eq:antipode}
\end{equation}
with constraint Jacobian $\det=2M^2|z_5|^2$, and the phase-space
measure reduces to
\begin{equation}
\dLIPS \;=\; \frac{\dd^2 z_5}{8\pi^2\,(1+|z_5|^2)^2},
\qquad \int \dLIPS=\frac{1}{8\pi}.
\label{eq:lips}
\end{equation}
\end{theorem}

\begin{remark}
Equation~\eqref{eq:lips} is the round (Fubini--Study) measure on the
celestial sphere: \emph{the cut is uniform on the sphere}, and the two
cut particles sit at exactly antipodal points.  The antipodal map ---
the geometric core of the shadow/CPT structure --- is thus an output of
momentum conservation, not an input.  All four components of the
constraint vanish identically on \eqref{eq:antipode}
(verified symbolically), and the total $\tfrac{1}{8\pi}$ is recovered
exactly.
\end{remark}

% =====================================================================
\section{Link 2: the Mellin image of the cut and the shadow locus}
\label{sec:link2}
% =====================================================================

Define the Mellin image of the cut measure, weighting the two cut
energies by boost weights,
\begin{equation}
\Phi(\Delta_5,\Delta_6;M)\;:=\;\int \dLIPS\;
\omega_5^{\Delta_5-1}\,\omega_6^{\Delta_6-1}.
\end{equation}

\begin{theorem}[Mellin image of the two-particle cut]
\label{thm:mellincut}
\begin{equation}
\Phi(\Delta_5,\Delta_6;M)
\;=\;
\frac{1}{8\pi}\,\Bigl(\frac{M}{2}\Bigr)^{\Delta_5+\Delta_6-2}\,
\frac{\Gamma(\Delta_5)\,\Gamma(\Delta_6)}{\Gamma(\Delta_5+\Delta_6)} .
\label{eq:phi}
\end{equation}
\end{theorem}

\begin{proof}
Insert \eqref{eq:antipode} into \eqref{eq:lips}; with $u=|z_5|^2$ the
angular integral is trivial and the radial integral is the Euler Beta
function $B(\Delta_6,\Delta_5)$.  Verified by quadrature at five points
with complex $\Delta$ on and off the principal series to relative error
$10^{-20}$--$10^{-31}$.
\end{proof}

\begin{corollary}[Cutkosky closure on the bubble]
\label{cor:bubble}
$\Phi(1,1;M)=\tfrac{1}{8\pi}=2\operatorname{Im}I_2(s)$ for the massless
bubble, whose absorptive part is the constant
$\operatorname{Im}I_2=\tfrac{1}{16\pi}$.
\end{corollary}

\begin{remark}[The true meaning of $\Delta_5+\Delta_6=2$]
\label{rem:shadowlocus}
Three statements replace the earlier ``residue at the shadow pole'':
\begin{enumerate}
\item The $M$-dependence of the cut's Mellin image is
$(M/2)^{\Delta_5+\Delta_6-2}$.  The locus $\Delta_5+\Delta_6=2$ is the
locus of \emph{exact scale invariance of the cut}; for the bubble,
whose absorptive part is constant in $s$, the physical cut sits
precisely there.  Nothing is localized by a pole; there is no residue
to take.
\item The pairing $(\Delta,2-\Delta)$ of the two cut legs is the shadow
pairing of the massless completeness relation on the principal series
\cite{LamShao2018,CelestialOptical2024}; the associated normalization
$N_{h,m=0}$ must be carried when the sewing is written in a
celestial-native operator form.
\item The $\Gamma\Gamma/\Gamma$ structure asserted for celestial
three-point coefficients in every earlier draft is \emph{derived} in
\eqref{eq:phi} from pure phase space.  On the shadow locus with
$\Delta_5=1+i\lambda$,
\begin{equation}
\Gamma(1+i\lambda)\Gamma(1-i\lambda)
=\frac{\pi\lambda}{\sinh(\pi\lambda)}\;=\;P(\lambda),
\label{eq:plambda}
\end{equation}
which identifies the spectral weight of the earlier drafts as the
phase-space/OPE factor of the cut --- see
Sec.~\ref{sec:relation}.
\end{enumerate}
\end{remark}

% =====================================================================
\section{Link 3: the scalar box cut and its Mellin poles}
\label{sec:link3}
% =====================================================================

\subsection{The regulated box and its $s$-channel cut}

Take the all-incoming planar box with massless external legs
$k_1,\dots,k_4$, massless propagators $\ell^2$ and $(\ell-k_1-k_2)^2$
(the cut lines), and mass-$\mu$ propagators $(\ell-k_1)^2-\mu^2$ and
$(\ell+k_4)^2-\mu^2$ (the tree lines).  The regulator keeps the cut
legs massless --- so Theorem~\ref{thm:measure} applies verbatim ---
while removing the collinear divergences of the strictly massless box
and rendering every comparison scheme-free.  Channel invariants:
$s=(k_1+k_2)^2>0$, and the second (massive-line) channel invariant is
$u=(k_2+k_3)^2<0$; physical $s$-channel kinematics has
$-s<t<0$, $s+t=-u$.  Define the pure cut integral
\begin{equation}
C(s,t,\mu^2)\;:=\;\int \dLIPS\;
\frac{1}{\bigl[(\ell_5-k_1)^2-\mu^2\bigr]\,
\bigl[(\ell_5+k_4)^2-\mu^2\bigr]}\;>0 .
\label{eq:cutdef}
\end{equation}

\begin{theorem}[Exact box cut]
\label{thm:boxcut}
With $c=\mu^2(s+\mu^2)$ and $d=s(s+t)=-su>0$,
\begin{equation}
C(s,t,\mu^2)
=\frac{1}{8\pi}\int_0^1\!\frac{\dd x}{\,c+d\,x(1-x)\,}
=\frac{1}{8\pi}\,\frac{4}{\sqrt{d(d+4c)}}\,
\operatorname{artanh}\!\sqrt{\frac{d}{d+4c}} ,
\label{eq:cutclosed}
\end{equation}
and in celestial variables (leg $1$ at $z_1=0$, leg $4$ at
$z_4=-\cot(\theta^*/2)$, $\omega_4=\tfrac{M}{2}\sin^2(\theta^*/2)$),
\begin{equation}
C=\int\!\frac{\dd^2 z}{8\pi^2(1+|z|^2)^2}\,
\frac{1}{\Bigl[\dfrac{M^2|z|^2}{1+|z|^2}+\mu^2\Bigr]
\Bigl[\,4\,\omega_5\omega_4|z-z_4|^2+\mu^2\Bigr]},
\qquad \omega_5=\frac{M}{2(1+|z|^2)} .
\label{eq:cutcelestial}
\end{equation}
The two representations and a direct center-of-mass angular integral
agree at four kinematic points to relative error
$10^{-23}$--$10^{-26}$.  As $\mu\to0$,
$C\to \log\!\bigl[s(s+t)/(\mu^2 s)\bigr]/\bigl[4\pi s(s+t)\bigr]$
(collinear logarithm; verified to $10^{-5}$ relative at
$\mu^2=10^{-5}$).
\end{theorem}

Equation~\eqref{eq:cutcelestial} is the honest ``shadow discontinuity''
of the program: a rational integrand on the celestial sphere, built
from tree propagators via \eqref{eq:qq}, integrated against the derived
measure \eqref{eq:lips}.  It contains no inserted factors.

\subsection{Mellin poles of the cut}

At fixed scattering angle, $\kappa:=(s+t)/s=(1+\cos\theta^*)/2$ is
constant, and we define
$\cM_C^{\angle}(\sigma)=\int_0^\infty\dd s\,s^{\sigma-1}C(s)$ on the
strip $0<\operatorname{Re}\sigma<2$.

\begin{theorem}[Universal double pole; infrared confinement to the
simple pole]
\label{thm:poles}
$\cM_C^{\angle}(\sigma)$ continues meromorphically to
$\sigma=2$ with
\begin{equation}
\cM_C^{\angle}(\sigma)
=\frac{1}{4\pi\kappa}
\left[\frac{1}{(2-\sigma)^2}
+\frac{\log(\kappa/\mu^2)}{2-\sigma}\right]
+\;\text{regular}.
\label{eq:poles}
\end{equation}
The double-pole coefficient $\tfrac{1}{4\pi\kappa}$ is
$\mu$-independent (universal); the collinear logarithm resides entirely
in the simple pole.  Verified to $7$ (double pole) and $8$ (simple
pole) significant digits at two kinematics, using an exact analytic
tail: naive numerical limits fail because at $\sigma=2-\epsilon$ the
Mellin mass sits at $s\sim e^{1/\epsilon}$.
\end{theorem}

The value of $\cM_C^{\angle}(1+i\lambda)$ on the principal line is
finite and smooth; it is the honest spectral density of the cut
(the object earlier drafts called $K_1$), now obtained from the true
side of the equation.

% =====================================================================
\section{Link 4: the loop from its cut, and the celestial dispersion
relation}
\label{sec:link4}
% =====================================================================

Write the box in Feynman-parametric form,
$I_4=\tfrac{i}{16\pi^2}J(s,u)$ with
\begin{equation}
J(s,u)=\int_{\Sigma_3}\!\dd^4x\;\delta\Bigl(1-\textstyle\sum x_i\Bigr)\,
\Delta(x)^{-2},\qquad
\Delta=\mu^2(x_2+x_4)-s\,x_1x_3-u\,x_2x_4 ,
\label{eq:parametric}
\end{equation}
real and positive in the Euclidean region $s,u<0$.

\begin{lemma}[Cutkosky normalization, derived and anchored]
\label{lem:cutkosky}
$\operatorname{Im}J(s,u)=8\pi^2\,C(s,u,\mu^2)$ for $s>0$, $u<0$.  The
constant follows from the cutting rule
$\prod_{\rm cut}(D_i+i\varepsilon)^{-1}\to(-2\pi i)\,\delta^+(D_i)$
together with the measure identity
$\int\frac{\dd^4\ell}{(2\pi)^4}(2\pi)^2\delta^+(D_1)\delta^+(D_3)f
=\int\dLIPS\,f$, and is anchored \emph{exactly} on the bubble:
$\operatorname{Im}J_2=\pi=8\pi^2\cdot\tfrac{1}{8\pi}$.
\end{lemma}

\begin{theorem}[Loop reconstruction; zero free parameters]
\label{thm:disp}
For $s<0$, $u<0$ (so that only the massless $s$-channel cut exists, the
massive channel opening only at $u\ge 4\mu^2$), the unsubtracted
fixed-$u$ dispersion relation holds:
\begin{equation}
J(s,u)\;=\;8\pi\int_0^{\infty}\!\dd s'\,
\frac{C(s',u,\mu^2)}{s'-s}.
\label{eq:disp}
\end{equation}
Verified against the direct parametric integral
\eqref{eq:parametric} at four Euclidean points to relative error
$0$--$1.6\times10^{-21}$, with the inner Feynman-parameter integral
done in the exact closed form
$\int_0^1\dd y\,[c+dy(1-y)]^{-2}
=\tfrac{8}{\sqrt d\,w^{3/2}}\operatorname{artanh}\sqrt{d/w}
+\tfrac{2}{cw}$, $w=d+4c$.
\end{theorem}

\begin{theorem}[Celestial dispersion relation]
\label{thm:celdisp}
Let $\cM_J(\sigma)=\int_0^\infty\dd S\,S^{\sigma-1}J(-S,u)$ and
$\cM_C(\sigma)=\int_0^\infty\dd s'\,s'^{\sigma-1}C(s',u,\mu^2)$ at
fixed $u<0$.  Then on the common strip $0<\operatorname{Re}\sigma<1$,
\begin{equation}
\boxed{\;\;
\cM_J(\sigma)\;=\;\frac{8\pi^2}{\sin(\pi\sigma)}\;\cM_C(\sigma).
\;\;}
\label{eq:celdisp}
\end{equation}
Verified to $10^{-11}$ (real $\sigma$) and $3\times10^{-7}$
(complex $\sigma$; nested-quadrature limited).
\end{theorem}

\begin{proof}
Mellin transform of \eqref{eq:disp} using the Stieltjes kernel
$\int_0^\infty\dd S\,S^{\sigma-1}(s'+S)^{-1}
=s'^{\sigma-1}\pi/\sin(\pi\sigma)$ and Fubini on the positive
integrand.
\end{proof}

\begin{remark}
Equation~\eqref{eq:celdisp} is the corrected content of the tree-to-loop
program: \emph{in Mellin space, the loop amplitude is the Mellin
transform of the tree-level cut multiplied by a universal elementary
kernel}.  The poles of $1/\sin(\pi\sigma)$ at integer $\sigma$ are the
familiar $\beta$-plane poles of celestial amplitudes; loop corrections
promote them to higher order via the pole structure of $\cM_C$
(Thm.~\ref{thm:poles}), consistent with the known analytic structure of
loop-level celestial amplitudes \cite{AHH2020,Chang2021}.
\end{remark}

% =====================================================================
\section{Relation to the shadow-discontinuity program and the
companion drafts}
\label{sec:relation}
% =====================================================================

\subsection{What survives of v14--v18}
\begin{itemize}
\item \textbf{Topology selection (survives, reinterpreted).}  The v18
statement that choosing a factorization channel of higher-point tree
data and closing a designated leg pair selects a definite loop topology
is correct, and is here realized concretely: the choice of cut channel
fixes the two uncut tree propagators in \eqref{eq:cutdef}; sewing the
pair over the derived phase space \eqref{eq:lips} closes the cycle.
The counting (tree with $4+2L$ legs $\to$ $L$-loop four-point topology
by $L$ pairwise closures) remains a plausible \emph{conjecture} beyond
one loop; no theorem is claimed here for $L\ge2$.
\item \textbf{The shadow pairing (survives, derived).}
$\Delta_6=2-\Delta_5$ is the completeness-relation pairing of the cut
pair; the antipodal map $z_6=-1/\bar z_5$ is forced by momentum
conservation (Thm.~\ref{thm:measure}).
\item \textbf{The locus $\Delta_5+\Delta_6=2$ (corrected).}  It is the
scale-invariance locus of the cut (Rem.~\ref{rem:shadowlocus}), not a
simple pole of the six-point tree from which a residue is extracted.
Drafts v14--v15, built on the residue mechanism and on the discredited
closed-form and \texttt{check\_box} benchmarks, are withdrawn.  The v18
separation of extraction/sewing/evaluation and its disavowal of the
mislabeled benchmarks are endorsed and completed here: the
``direct-shadow, no-Feynman-subroutine'' benchmark that v18 called for
is precisely the celestial integral \eqref{eq:cutcelestial}, now
implemented and verified.
\end{itemize}

\subsection{The ``Plancherel loop measure'' (haar\_qg v2.1.5.1)}
The spectral weight $P(\lambda)=\pi\lambda/\sinh(\pi\lambda)$ is
\emph{real and derived}: it is
$\Gamma(1+i\lambda)\Gamma(1-i\lambda)$, the phase-space/OPE factor of
the cut on the shadow locus, Eq.~\eqref{eq:plambda}.  It is
\emph{not} the $\mathrm{SL}(2,\C)$ Plancherel density (which is
$\propto\lambda^2$ for the scalar principal series), it is not a
replacement for the Feynman loop measure, and it does not render loop
expansions ultraviolet finite: the reconstruction
\eqref{eq:disp}--\eqref{eq:celdisp} reproduces the standard loop
integral exactly, divergences and all (here regulated by $\mu$).  The
claims $\cM_L\le(1/8)^L$ ``UV finiteness with no counterterms,'' the
identification of $\cM_1=1/8,\ \cM_2=1/90$ as physical loop
normalizations, and the fifteen-digit box verification (which reduces
to \texttt{check\_box}) are withdrawn.  The iterated moments of
$P(\lambda)$ and their Bernoulli-fraction values remain correct as
special-function identities.

\subsection{The kinematic block (v1.1) and the blackbody draft}
The special-function content of the kinematic-block paper --- the
reduction of the chiral block $k_h(z)$ on the principal series to the
conical Legendre function $Q_{h-1}$, the shadow involution as the
degree symmetry $\nu\mapsto-\nu-1$, and the Mehler--Fock structure of
the shadow-symmetric part --- is correct classical mathematics and is
retained: it is exactly the toolkit required for the celestial-native
partial-wave decomposition of the cut \eqref{eq:cutcelestial}
(companion work; the genuine $\mathrm{SL}(2,\C)$ Plancherel density
enters there, resolving the long-standing ``missing normalization'').
The blackbody draft's identity relating $P(\lambda)$ to a modular
(Bisognano--Wichmann, $\beta=2\pi$) oscillator remains a correct
identity about the function \eqref{eq:plambda}; its framing as a cure
of ultraviolet divergences is withdrawn with the parent claim.

% =====================================================================
\section{What is not claimed}
\label{sec:notclaimed}
% =====================================================================
\begin{enumerate}
\item Results are proven for the massless bubble and the internally
$\mu$-regulated scalar box: one complete, fully verified worked case.
A general theorem for arbitrary topologies, masses, and multiplicities
is not claimed; the machinery (derived cut measure, closed
cut forms, exact-tail Mellin continuation, dispersion with derived
normalization) is expected to generalize and is left to future work.
\item The strictly massless box requires dimensional or mass
regularization of its collinear cut; Theorem~\ref{thm:poles} shows
precisely how the infrared logarithm is confined to the subleading
Mellin pole, but no scheme-independence statement beyond that is made.
\item The multi-loop pattern ($(4+2L)$-point tree data $\to$ $L$-loop
integrands) is a conjecture.
\item No connection to the Riemann Hypothesis or to any Millennium
Prize problem is asserted in this paper.
\end{enumerate}

% =====================================================================
\section{Data and code availability}
\label{sec:code}
% =====================================================================
All results are reproduced by five self-contained Python scripts
(\texttt{mpmath}/\texttt{sympy}), archived with this paper and in
\texttt{GppVerify/discovery/shadow\_ope/}:
\texttt{celestial\_cut\_step1.py} (Thms.~\ref{thm:measure},
\ref{thm:mellincut}),
\texttt{box\_cut\_step2a.py} (Thm.~\ref{thm:boxcut}),
\texttt{mellin\_cut\_step2c.py} (Thm.~\ref{thm:poles}),
\texttt{dispersion\_step3a.py} (Thms.~\ref{thm:disp},
\ref{thm:celdisp}),
plus session notes \texttt{REDO\_NOTES\_LOOPS\_FROM\_TREES.md}.
The archived routines \texttt{box\_shadow\_1d}/\texttt{check\_box} are
deprecated and must not be cited as verification.

\subsection*{Acknowledgments}
This version was produced in an audited human--AI research session in
which all prior numerical claims of the series were independently
recomputed before any new claim was made.

\begin{thebibliography}{9}
\bibitem{Bern1994} Z.~Bern, L.~Dixon, D.~Dunbar, D.~Kosower,
``One-loop $n$-point gauge theory amplitudes, unitarity and collinear
limits,'' Nucl.\ Phys.\ B435 (1995) 59.
\bibitem{Pasterski2017} S.~Pasterski, S.-H.~Shao, A.~Strominger,
``Flat space amplitudes and conformal symmetry of the celestial
sphere,'' Phys.\ Rev.\ D96 (2017) 065026.
\bibitem{Raclariu2021} A.-M.~Raclariu, ``Lectures on celestial
holography,'' arXiv:2107.02075.
\bibitem{LamShao2018} H.~T.~Lam, S.-H.~Shao, ``Conformal basis,
optical theorem, and the bulk point singularity,'' Phys.\ Rev.\ D98
(2018) 025020 [arXiv:1711.06138].
\bibitem{CelestialOptical2024} ``The celestial optical theorem,''
arXiv:2404.18898.
\bibitem{AHH2020} N.~Arkani-Hamed, M.~Pate, A.-M.~Raclariu,
A.~Strominger (and related), on the analytic structure of celestial
amplitudes in the $\beta$-plane.
\bibitem{Chang2021} C.-M.~Chang et al., celestial dispersion
relations, 2021.
\end{thebibliography}

\end{document}
